\begin{table}[htbp]
\centering
\caption{Robustness of Main TVA Effect on Manufacturing Share}
\label{tab:robustness}
\begin{tabular}{lccc}
\toprule
Specification & Coefficient & SE & N \\
\midrule
Main DiD (all counties) & 0.0134 & (0.0073) & 5,291 \\
Border counties only & 0.0057 & (0.0108) & 1,005 \\
Donut (exclude 100-200km) & 0.0106 & (0.0065) & 4,637 \\
Pre-trend placebo (1920-1930) & 0.0027 & (0.0056) & 3,504 \\
\midrule
Randomization inference p-value & \multicolumn{3}{c}{0.002} \\
Wild cluster bootstrap p-value & \multicolumn{3}{c}{0.006} \\
\bottomrule
\end{tabular}
\begin{minipage}{0.95\textwidth}
\vspace{0.5em}
\footnotesize
\textit{Notes:} All specifications include county and year fixed effects with state-clustered standard errors. The dependent variable is manufacturing employment share. Row 1 is the baseline from Table \ref{tab:main_did}. Border counties restricts to TVA core and adjacent non-TVA counties. The donut excludes counties in the 100--200 km ambiguous zone. The pre-trend placebo uses only 1920 and 1930 (both pre-TVA). Randomization inference permutes TVA assignment across counties (500 permutations). Wild cluster bootstrap uses Rademacher weights at the state level (999 replications).
\end{minipage}
\end{table}
